\documentclass[11pt,a4paper]{article}\usepackage[utf8]{inputenc}\usepackage{amsmath, amssymb, amsfonts}\usepackage{geometry}\usepackage{graphicx}\usepackage{hyperref}\geometry{margin=1in}\title{High-Dimensional Effects in Machine Learning: \ From Random Walks to Hardware Bias}\author{Research Educational Series}\date{February 2026}\begin{document}\maketitle\section{Introduction}In modern machine learning, embedding spaces typically operate in dimensions ($H$) ranging from 1,024 to 8,192. At these scales, geometry behaves in ways that contradict our 3D physical intuitions. This article provides a framework for distinguishing ``real structure'' from finite-sample noise using high-dimensional analysis.\section{The Mathematical Baseline: The Noise Floor}To evaluate a model, we must first define the ``null hypothesis''---what the data looks like if it contains no meaningful structure.\subsection{The Zero-Mean Noise Assumption}If tokens $x_1, x_2, \dots, x_N$ (where $x_i \in \mathbb{R}^H$) are pure noise, they are considered isotropic, meaning they have no preferred direction. Adding these vectors together results in a \textbf{Random Walk} in $H$-dimensional space.\subsection{Scaling of the Mean Vector}Define the batch mean $\mu$ as:\begin{equation}\mu = \frac{1}{N} \sum_{i=1}^{N} x_i\end{equation}For a random walk of $N$ steps, the expected magnitude of the mean vector does not vanish to zero but scales by $\sqrt{1/N}$. The expected norm is defined as:\begin{equation}|\mu|_2 \sim \sqrt{\frac{\mathbb{E}|x|_2^2}{N}}\end{equation}\section{The $R$ Metric: Signal-to-Noise Ratio}We quantify ``structured bias'' by comparing the observed mean to the expected noise floor:\begin{equation}R = \frac{|\mu|_2}{\sqrt{\mathbb{E}|x|_2^2 / N}}\end{equation}\begin{itemize}\item \textbf{$R \approx 1$}: The mean is at the noise level; no coherent structure exists.\item \textbf{$R \gg 1$}: The mean is far above the noise floor, indicating a \textbf{coherent structured bias}.\end{itemize}\section{The Role of Dimensionality ($H$)}The dimensionality $H$ fundamentally alters the reliability of these metrics.\subsection{Concentration of Measure}In high dimensions, random vectors become nearly orthogonal ($90^\circ$) with a probability near 1. This leads to efficient noise cancellation.\begin{itemize}\item Individual vector norm: $\|x\|_2 \sim \sqrt{H}$.\item Expected mean norm: $\|\mu\|_2 \sim \sqrt{H/N}$.\end{itemize}\section{Impact of Hardware Bias in AI Silicon}In hardware-aware optimization (e.g., $FP4$ or $FP8$), micro-biases ($\epsilon$) in rounding or casting accumulate.\subsection{Parallel Accumulation}A bias $\epsilon$ applied to all $H$ dimensions simultaneously creates a total Euclidean shift:\begin{equation}\text{Shift} = \sqrt{H} \cdot \epsilon\end{equation}For $H=4096$ and $\epsilon=0.0001$, the shift is $\sqrt{4096} \cdot 0.0001 = 0.0064$. Over many layers, this accumulation leads to \textbf{Dimensional Collapse}.\section{Conclusion}The $R$ ratio serves as a vital ``health check'' for high-dimensional models. It allows researchers to detect when hardware constraints introduce biases that interfere with the model's ability to represent the physical world.\end{document}